% Part: first-order-logic
% Chapter: syntax-and-semantics
% Section: first-order-languages

\documentclass[../../../include/open-logic-section]{subfiles}

\begin{document}

\olfileid{fol}{syn}{fol}

\olsection{Talen van de logica van de eerste orde}


Om uitdrukkingen van de logica van de eerste orde te maken, beginnen we
met een basisvoorraad tekens. Die bevat \emph{!!{variable}s},
\emph{!!{constant}s}, \emph{!!{predicate}s} en soms
\emph{!!{function}s}. We voegen daar logische verbindingstekens,
kwantoren en leestekens zoals haakjes en komma's aan toe. Zo bouwen we
\emph{termen} en \emph{!!{formula}s} op.

\begin{explain}
Los gezegd staan !!{predicate}s voor eigenschappen en relaties. Een
!!{constant} is een naamteken voor één object, en een !!{function} is
een teken voor een functie. Al deze tekens, behalve de
!!{identity}~$\eq$, heten \emph{niet-logische symbolen}. Samen vormen ze
een taal. Welke eerste-ordetaal~$\Lang L$ we hebben, wordt volledig
bepaald door haar niet-logische symbolen. In het algemeenste geval
bevat $\Lang L$ van iedere soort oneindig veel symbolen.
\end{explain}

In het algemene geval gebruiken we voor de logica van de eerste orde
de volgende tekens:

\begin{enumerate}
\item Logische tekens
\begin{enumerate}
\item Logische verbindingstekens:
  \startycommalist
  \iftag{prvNot}{\ycomma $\lnot$ (negatie)}{}%
  \iftag{prvAnd}{\ycomma $\land$ (conjunctie)}{}%
  \iftag{prvOr}{\ycomma $\lor$ (disjunctie)}{}%
  \iftag{prvIf}{\ycomma $\lif$ (!!{conditional})}{}%
  \iftag{prvIff}{\ycomma $\liff$ (!!{biconditional})}{}%
  \iftag{prvAll}{\ycomma $\lforall$ ('voor alle'-kwantor)}{}%
  \iftag{prvEx}{\ycomma $\lexists$ ('er bestaat'-kwantor)}{}.
\tagitem{prvFalse}{Het vaste propositionele teken voor !!{falsity}~$\lfalse$.}{}
\tagitem{prvTrue}{Het vaste propositionele teken voor !!{truth}~$\ltrue$.}{}
\item Het tweeplaatsige teken voor !!{identity}~$\eq$.
\item Een !!{denumerable}s verzameling met !!{variable}s: $\Obj v_0$, $\Obj v_1$, $\Obj
  v_2$, \dots
\end{enumerate}
\item Niet-logische symbolen. Samen vormen die de \emph{standaardtaal}
  van de logica van de eerste orde
\begin{enumerate}
\item Een !!{denumerable}s verzameling !!{predicate}s met $n$ argumentplaatsen voor ieder $n>0$: $\Obj
  A^n_0$, $\Obj A^n_1$, $\Obj A^n_2$, \dots
\item Een !!{denumerable}s verzameling met !!{constant}s: $\Obj c_0$, $\Obj c_1$, $\Obj
  c_2$, \dots.
\item Een !!{denumerable}s verzameling !!{function}s met $n$ argumentplaatsen voor ieder $n>0$:
  $\Obj f^n_0$, $\Obj f^n_1$, $\Obj f^n_2$, \dots
\end{enumerate}
\item Leestekens: (, ) en de komma.
\end{enumerate}

De meeste definities en resultaten geven we voor de volledige
standaardtaal van de logica van de eerste orde. Voor een bepaalde
toepassing kunnen we ook een kleinere taal nemen, met maar enkele
!!{predicate}s, !!{constant}s en !!{function}s.

\begin{ex}
De taal~$\Lang L_A$ van de rekenkunde heeft één !!{predicate} met twee
argumentplaatsen, namelijk~$<$. Ze heeft ook één !!{constant}~$\Obj 0$,
één !!{function} met één argumentplaats~$\prime$ en twee !!{function}s
met twee argumentplaatsen, $+$ en~$\times$.
\end{ex}

\begin{ex}
De taal van de verzamelingenleer~$\Lang L_Z$ bevat maar één
!!{predicate} met twee argumentplaatsen:~$\in$.
\end{ex}

\begin{ex}
De taal van ordeningen~$\Lang L_\le$ bevat maar één !!{predicate} met
twee argumentplaatsen:~$\le$.
\end{ex}

Ook deze schrijfwijzen zijn afspraken. In de officiële notatie zijn
$<$, $\in$ en $\le$ alleen andere namen voor $\Obj A^2_0$. Zo staat
$\Obj 0$ voor $\Obj c_0$, $\prime$ voor $\Obj f^1_0$, $+$ voor
$\Obj f^2_0$ en $\times$ voor $\Obj f^2_1$.

\iftag{defNot,defOr,defAnd,defIf,defIff,defTrue,defFalse,defEx,defAll}{%
Naast de basisverbindingstekens en
\iftag{notprvEx,notprvAll}{kwantor}{kwantoren} hierboven gebruiken we
ook de volgende \emph{gedefinieerde} symbolen:
\startycommalist
  \iftag{defNot}{\ycomma $\lnot$ (negatie)}{}%
  \iftag{defAnd}{\ycomma $\land$ (conjunctie)}{}%
  \iftag{defOr}{\ycomma $\lor$ (disjunctie)}{}%
  \iftag{defIf}{\ycomma $\lif$ (!!{conditional})}{}%
  \iftag{defIff}{\ycomma $\liff$ (!!{biconditional})}{}%
  \iftag{defAll}{\ycomma $\lforall$ ('voor alle'-kwantor)}{}%
  \iftag{defEx}{\ycomma $\lexists$ ('er bestaat'-kwantor)}{}%
  \iftag{defFalse}{\ycomma !!{falsity}~$\lfalse$}{}%
  \iftag{defTrue}{\ycomma !!{truth}~$\ltrue$}}{}.

\begin{tagblock}{defNot,defOr,defAnd,defIf,defIff,defTrue,defFalse,defEx,defAll}
\begin{explain}
Een gedefinieerd symbool hoort officieel niet bij de taal. We voeren
het in als afkorting. Daarmee kunnen we formules korter opschrijven die
met alleen de basissymbolen erg lang zouden worden. Nog belangrijker:
ook bewijzen worden korter. Voor ieder basissymbool is in een bewijs
namelijk een apart geval nodig. Meer basissymbolen leveren dus langere
bewijzen op.
\end{explain}
\end{tagblock}

% Alternate symbols

\begin{intro}
Misschien ken je andere termen en symbolen. Boeken en docenten
gebruiken voor ``negatie'' vaak $\sim$, $\neg$ of~!\ en voor
``conjunctie'' vaak $\wedge$, $\cdot$ of $\&$. Voor een ``als-dan-teken''
of ``implicatie'' zijn $\rightarrow$, $\Rightarrow$ en $\supset$
gebruikelijk.
\iftag{prvIff,defIff}{Voor ``biconditionele verbinding'', ``bi-implicatie''
  of ``(materiële) equivalentie'' gebruikt men $\leftrightarrow$,
  $\Leftrightarrow$ en $\equiv$.}{}
\iftag{prvFalse,defFalse}{Het teken $\lfalse$ heet ook wel
  ``onwaarheid'', ``falsum'', ``absurdum'' of ``bottom''.}{}
\iftag{prvTrue,defTrue}{Het teken $\ltrue$ heet ook wel
  ``waarheid'', ``verum'' of ``top''.}{}

Voor !!{constant}s, die soms namen heten, gebruikt men meestal kleine
letters aan het begin van het Latijnse alfabet, zoals $a$, $b$ en $c$.
Voor !!{variable}s gebruikt men meestal letters aan het eind, zoals
$x$, $y$ en $z$. Een kwantor hoort bij een !!{variable}, bijvoorbeeld
$x$. De 'voor alle'-kwantor kan worden geschreven als $\forall x$,
$(\forall x)$, $(x)$, $\Pi x$ of $\bigwedge_x$. Voor de 'er
bestaat'-kwantor komen $\exists x$, $(\exists x)$, $(Ex)$, $\Sigma x$
en $\bigvee_x$ voor.
\end{intro}

\begin{explain}
We kunnen alle propositionele operatoren en beide kwantoren als
basistekens nemen. We kunnen ook minder basistekens kiezen en de andere
!!{operator}s daarmee definiëren. Een verzameling Booleaanse operatoren
heet ``functioneel volledig'' als je er iedere Booleaanse formule mee
kunt uitdrukken. Voorbeelden zijn $\{ \lnot, \lor \}$,
$\{ \lnot, \land \}$ en $\{ \lnot, \lif\}$. Combineer zo'n verzameling
met een van beide kwantoren, dan krijg je een eerste-ordetaal met
volledige uitdrukkingskracht.

Misschien ken je nog twee !!{operator}s: de Sheffer-streep~$|$, genoemd
naar Henry Sheffer, en de pijl van Peirce~$\downarrow$, die ook de dolk
van Quine heet. Hun gebruikelijke betekenissen zijn respectievelijk
``nand'' en ``nor''. Met elk van deze operatoren afzonderlijk kun je
alle Booleaanse formules uitdrukken: ieder is dus functioneel volledig.
\end{explain}
\end{document}
